\section{Related Work}
\label{sec:related}

\subsection{Reuse and adaptation of agent skills}

Gao et al.~\cite{gao2026registry} present the first study of skills as engineered
artifacts. They collected 18,463 skills from the skills.sh registry and 23,199 skills from
5,876 GitHub repositories, and linked a repository skill to a registry skill when both
carried the same YAML \texttt{name}, using a longest-common-subsequence similarity to choose
among same-name candidates and to reject pairs below 0.1. This recovered 3,709 reuse links.
Most reused skills turned out to be near-verbatim copies, 53\% were never modified after
adoption, and among reused skills that were never updated locally, 40.2\% had a registry
source that changed afterwards, changes the copies never received. The linkage only sees a
copy whose source is published on skills.sh and that kept the source's name. The authors
acknowledge that it misses renamed and heavily edited skills and assumes that reuse flows
from the registry into repositories; 76\% of the repository skills they collected stayed
unlinked, and some of those may be undetected reuse.

Wu et al.~\cite{wu2026adaptation} study how downstream developers adapt published skills,
using public forks of six popular skill repositories as evidence. From 1,126 adaptation
instances they build a taxonomy of 46 adaptation patterns in 13 families, and they report
that 18.6\% of adaptations add security-sensitive content to the instruction text. They call
for tools that detect inconsistent modifications. Because the method relies on GitHub's fork
relation, it observes only copies made by forking those six repositories. A skill pasted into
an unrelated repository leaves no fork record.

\subsection{Clone detection for agent skills}

Zhu et al.~\cite{zhu2026skillreuse} propose SkillReuse, the work closest to ours. It
represents a skill's YAML metadata, prose, and code separately, combines the per-channel
similarities with a flat TF-IDF similarity and structural features in a logistic regression,
and reaches an F1 of 0.939 on SkillReuse-Bench, a set of 300 labelled pairs. Applied to
137,470 skills crawled from the SkillsMP marketplace, it places 66.8\% of skills in at least
one clone pair, most of them across authors. The authors are explicit about what clone
detection does not establish. Structural similarity is not derivation, so detected pairs may
include convergent development; the clone relation is undirected and does not say which skill
came first; and the analysis uses the \texttt{SKILL.md} file alone, without repository
history. Two properties of the benchmark also matter for our problem. Half of its 150
positive pairs are synthetic mutations rather than real edits, and its hard negatives are
same-category pairs whose similarity lies between 0.15 and 0.45. Highly similar skills that
were written independently, such as two skills generated from one template, are therefore
not represented. Finally, families formed as connected components of the clone graph chain
together: the largest one contains 61,651 skills.

\subsection{Code clones and reuse on GitHub}

Clone detection for source code is a mature field. Roy et al.~\cite{roy2009comparison}
define the standard clone types and compare detection techniques, and SourcererCC scales
token-based detection of near-miss clones to large code corpora~\cite{sajnani2016sourcerercc}.
Lopes et al.~\cite{lopes2017dejavu} showed that file-level duplication is pervasive on GitHub.
Kim et al.~\cite{kim2005genealogy} introduced clone genealogies, which follow groups of
clones across versions and separate consistent changes, applied to every member of a group,
from inconsistent ones. These techniques were designed for source code, where tokens and
syntax carry most of the evidence. Skills are mostly natural language and are often written
from shared scaffolds, which makes surface similarity a weaker signal of shared origin.
